\documentclass[10pt,a4paper]{amsart}
\usepackage[margin=19mm]{geometry}
\usepackage{amsmath,amssymb,mathtools}
\usepackage[scr=rsfs,cal=euler]{mathalfa}
\usepackage{xfrac}
\usepackage[unicode,hidelinks,bookmarks=false]{hyperref}
\newcommand{\ie}{i.e.\ }
\newcommand{\cf}{cf.\ }
\newcommand{\resp}{respectively\ }
\newcommand{\Subsection}{\subsection}
\providecommand{\nobreakdash}{\nobreak\hbox{-}}
\setlength{\emergencystretch}{2em}
\multlinegap0pt
\usepackage{amssymb}
\usepackage{mathtools}
\newtheorem{proposition}{Proposition}
\newtheorem{definition}{Definition}
\title{A transform approach to the supercooled Stefan problem}
\author{Thomas Blore, Ben Hambly}
\date{Source: arXiv:2609.24297v1 (CC BY 4.0)}
\begin{document}
\maketitle

\noindent\emph{Source and licence.} A transform approach to the supercooled Stefan problem. Version arXiv:2609.24297v1 (CC BY 4.0), distributed by arXiv under the Creative Commons Attribution 4.0 International licence, http://creativecommons.org/licenses/by/4.0/, as recorded in the arXiv licence metadata for this exact version and shown on the abstract page https://arxiv.org/abs/2609.24297v1. Full licence notice: \texttt{licenses/arxiv-2609.24297-CC-BY-4.0.txt}.\par\smallskip

\noindent\textit{Editorial scope.} Counted unit: the article's proposition Ito (source0.tex lines 121--126). The article's complete original proof of it is delivered with it (source0.tex lines 339--374, one \texttt{proof} environment of 36 lines, which the article places away from the statement). No mathematics is written, repaired or re-derived: the statement and the proof are the article's own lines, in the article's own notation, and no retained line is substituted or re-flowed.  Retained uncounted as the source-local context the proof needs: lines 82--87; lines 97--101.  Nothing was omitted from any retained span, so the package carries no omission record.  No retained line contains a citation, so the wrapper carries no bibliography and no bibliography entry.  No retained line contains a figure environment, an image-inclusion command or drawing code, so no image is delivered and the package declares no asset dependency.  Editorial formatting only, in addition: an AMS \texttt{amsart} preamble that restates the article's own declarations and macros used by the retained lines, namely the environments \texttt{proposition}, \texttt{definition} on separate counters, together with the standard packages those lines need.  The retained environments print in this document as Proposition 1, Definition 1 in order of appearance, whereas the article prints the same items with the numbering of its own full document.  The complete native source of the article as distributed in the arXiv e-print for this exact version is bundled unchanged as \texttt{sources/211175/source0.tex}.
\begin{align}
    \partial_t V(t,x) &=\frac{1}{2}\partial_{xx} V(t,x), \quad x>\Lambda_t,\nonumber\\
     \dot{\Lambda}_t &=\frac{1}{2}\partial_x V(t,\Lambda_t),\label{pde}\\
    V(t,\Lambda_t) &=0, \qquad\qquad t\geq 0, &\nonumber\\
    V(0,x) &=g(x),\qquad\;\; x>0.\nonumber
\end{align}

\begin{align}
    X^x_t &=x+B_t-\Lambda_t,\nonumber\\
    \tau^x &=\inf\{t\geq 0:X^x_t\leq 0\},\label{alternativetomvr}\\
    \Lambda_t &=\int_0^\infty \mathbb{P}(\tau^x\leq t)g(x)\textrm{d}x.\nonumber
\end{align}

\begin{proposition}[Laplace transform of solutions]
\label{Ito}
For all $\lambda>0$, any solution to \eqref{alternativetomvr} on the interval $[0,t]$ satisfies the following equation:
\begin{align}&\int_{\Lambda_t}^\infty V(t,x)e^{-\lambda x-\frac{\lambda^2t}{2}}\textrm{d}x+\int_0^\infty (1-g(x))e^{-\lambda x}\textrm{d}x-\frac{1}{\lambda}e^{-\lambda \Lambda_t-\frac{\lambda^2t}{2}}\\&=\frac{\lambda}{2}\int_0^te^{-\lambda\Lambda_{s^-}-\frac{\lambda^2s}{2}}\textrm{d}s+\sum_{s\leq t}e^{-\frac{\lambda^2s}{2}}\left(-\int_{\Lambda_{s^-}}^{\Lambda_s} V(s^-,x)e^{-\lambda x}\textrm{d}x-\frac{1}{\lambda}\Delta e^{-\lambda\Lambda_s}\right).\nonumber
\end{align}
\end{proposition}

\begin{proof}
We shall consider $C=1$ to reduce the notation. The proof is identical to that for $C\neq 1$. 
We shall consider $1/2<\alpha<1$ since the existence of similarity solutions ($\alpha=1/2$) and travelling wave solutions ($\alpha=1$) are well known. We shall construct the solution using the formula given by Proposition \ref{Ito}.

We first recall the definitions of a completely monotone function and a Bernstein function.
\begin{definition}
    A function $f$ is called completely monotone if it is continuous on $[0,\infty),$ infinitely differentiable on $(0,\infty)$, and satisfies $(-1)^n\frac{\textrm{d}^nf}{\textrm{d}x^n}(x)\geq 0$ for all non-negative integers $n$.
\end{definition}
\begin{definition}
    A Bernstein function is a non-negative differentiable function, with derivative which is completely monotone.
\end{definition}

We define a function $F:\mathbb{R}^+\times \mathbb{R^+}$ by
\begin{align*}
F(\lambda,t)&=e^{\lambda t^\alpha+\lambda^2t/2}\int_t^\infty e^{-\lambda s^\alpha-\lambda^2s/2}\alpha s^{\alpha-1}\textrm{d}s\\
    &=\frac{1}{\lambda}\int_0^\infty e^{-u}e^{-\frac{\lambda^2}{2}((\frac{u}{\lambda}+t^\alpha)^{1/\alpha}-t)} \textrm{d}u \textrm{ taking }u=\lambda(s^\alpha-t^\alpha).
\end{align*}

We claim that for each $t\geq 0$, $F(\lambda,t)$ is a completely monotone function of $\lambda$. To show this, it is sufficient to check that for all $u>0$, $t\geq 0$, \\$m( \lambda)=\frac{\lambda^2}{2}((\frac{u}{\lambda}+t^\alpha)^{1/\alpha}-t)$ is a Bernstein function in $\lambda$. It would then follow that
\\$\exp\left(-\frac{\lambda^2}{2}((\frac{u}{\lambda}+t^\alpha)^{1/\alpha}-t)\right)$ is completely monotone. Therefore the integrand in the above integral will also be completely monotone, and thus so is $F(.,t)$. The first derivative of $m(\lambda)$ can be explicitly computed, as can its inverse Laplace transform, from which it is easily seen that $m'(\lambda)$ is the Laplace transform of a positive function and hence a completely monotone function. It is then simple to see that $m(\lambda)$ is a Bernstein function.


We have shown that for all $t\geq 0$, $F(\lambda,t)$ is completely monotone in $\lambda$. Therefore, by Bernstein's theorem, for all $t\geq 0$, there exists a positive function $V_t:\mathbb{R}^+\rightarrow \mathbb{R^+}$ such that $V_t$ has the Laplace transform $F(\lambda,t)$. 

We define $g=V_0$. We write $(f,g)$ for the integral $\int_0^\infty f(x)g(x)dx$.  
By definition, the functions $V_t$ are such that for all $\lambda\geq 0$, $t\geq 0$:
\begin{align*}
    (V_t,e^{-\lambda x})&=(g,e^{-\lambda x})+\int_0^t\frac{1}{2}(V_s,\partial_{xx}\left(e^{-\lambda x}\right))\textrm{d}s\\&-\int_0^t\alpha s^{\alpha-1}(V_s,\partial_x\left(e^{-\lambda x}\right))\textrm{d}s-\int_0^t\alpha s^{\alpha -1}(e^{-\lambda x}|_{x=0})\textrm{d}s.
\end{align*}

By a standard density argument we extend this to 
\begin{align*}
    (V_t,\phi)=(g,\phi)+\int_0^t\frac{1}{2}(V_s,\phi'')\textrm{d}s-\int_0^t\alpha s^{\alpha-1}(V_s,\partial_x\phi)\textrm{d}s-\int_0^t\alpha s^{\alpha -1}\phi(0)\textrm{d}s
\end{align*}
for any $t\geq 0$ and $\phi$ a Schwartz function. Therefore, the function $V(t,x)=V_t(x+t^\alpha)$ and the free boundary $t^\alpha$, form a weak solution to the supercooled Stefan problem \eqref{pde} with initial data $g$. Since the function $t^\alpha$ is differentiable on $(0,\infty)$, it is straightforward to argue that $V(t,x)$ is a classical solution to the supercooled Stefan problem for $t>0$.
\end{proof}

\end{document}
